Throughout this paper, we adopt the following conventions. Let $\F$ denote an algebraically closed field and let $\operatorname{char}\F$ denote the characteristic of~$\F$. Let $\Z$ denote the set of integers and let $\N$ denote the set of nonnegative integers. The bracket $[\,,\,]$ stands for the commutator.

In this paper we consider the Racah algebra $\Re$ over~$\F$ defined by generators and relations as follows. The generators are $A$, $B$, $C$, $D$ and the relations assert that
 \begin{gather*}
 [A,B]=[B,C]=[C,A]=2D
\end{gather*}
 and that each of
\begin{gather*}
[A,D]+AC-BA, \qquad [B,D]+BA-CB, \qquad [C,D]+CB-AC
\end{gather*}
is central in $\Re$. The Racah algebra $\Re$ is a universal analog of the original Racah algebras which first appeared in~\cite{Levy1965}. In that paper, the original Racah algebras were used to establish a link between representation theory and the quantum mechanical coupling of three angular momenta. Since that time, the connections between the Racah algebras and many other areas have been explored. We mention a few of them here. Their connections with the additive double-affine Hecke algebra of type $\big(C_1^\vee,C_1\big)$, the Bannai--Ito algebra, and the Lie algebras $\mathfrak{su}(2)$, $\mathfrak{su}(1,1)$ were investigated in \cite{gvz2013,R&BI2015,zhedanov1988,Huang:R<BImodules}. Their realizations via the Racah polynomials, the intermediate Casimir operators, and the superintegrable models in two dimensions were presented in \cite{Galbert,R&LD2014,gvz2014,R&BI2015,quadratic1991}. For information concerning the higher rank Racah algebras, see~\cite{HR2017,HR2019}.
